\begin{abstract}

Large Language Models (LLMs) have recently gained significant attention
in financial forecasting because they can process both numerical and
textual information simultaneously. To improve reliability, researchers
have integrated retrieval techniques with LLMs, but the individual
contribution of each component remains unclear. This study addresses
this limitation by comparing four forecasting approaches under identical
experimental conditions: an LSTM model, a standalone LLM, an LLM
enhanced through Retrieval-Augmented Generation (RAG), and a hybrid
model combining the LSTM and the grounded LLM.

The evaluation was conducted using 300 test samples collected from five
U.S. stocks between 2019 and 2024. The hybrid model achieved the highest
directional compared with an always-up baseline models. However, statistical tests showed that the observed
improvements were not statistically significant.

Additional experiments examined hallucination behaviour, prediction
stability, and response quality. No fabricated values or false citations
were observed in either language model. Nevertheless, the introduction
of retrieval increased numerical unit errors from 2.6\% to 14.2\%.
Repeated experiments also revealed that the LLM produced inconsistent
predictions in a small number of cases, whereas the LSTM remained
completely deterministic.

The results demonstrate that the primary advantage of combining
retrieval and hybrid modelling lies in producing more reliable and
informative responses rather than substantially improving prediction
accuracy. The main contribution of this work is the introduction of a
reproducible framework that evaluates the individual impact of language
reasoning, retrieval grounding, and model fusion within financial
forecasting systems.

\end{abstract}
